\documentclass[10pt,a4paper]{amsart}
\usepackage{amsmath,amssymb,amsthm,graphicx}
% Retained minimal preamble for the Stacks Project chapter more-algebra.tex (More on Algebra).
% The chapter file pinned-original-source.tex carries only \input{preamble} on its first line, so no
% preamble text lives inside the chapter itself.  The Stacks Project uses one single shared preamble for
% every chapter; the snapshot stages it identically as sources/stacks-chow/preamble.tex and
% sources/stacks-etale-cohomology/preamble.tex.  The block below is copied mechanically, line for line,
% from sources/stacks-chow/preamble.tex lines 14-18 (the Xy-pic setup) and 154-217 (the native theorem
% declarations together with the mathematical macros).  Nothing was added, removed or redefined; the
% wrapper class and the amsmath/amssymb/amsthm/graphicx base packages are supplied by the assembler.
% Chapter: more-algebra.tex -- More on Algebra.
% For commutative diagrams we use Xy-pic
\usepackage[all]{xy}

% We use 2cell for 2-commutative diagrams.
\xyoption{2cell}

% Theorem environments.
%
\theoremstyle{plain}
\newtheorem{theorem}[subsection]{Theorem}
\newtheorem{proposition}[subsection]{Proposition}
\newtheorem{lemma}[subsection]{Lemma}

\theoremstyle{definition}
\newtheorem{definition}[subsection]{Definition}
\newtheorem{example}[subsection]{Example}
\newtheorem{exercise}[subsection]{Exercise}
\newtheorem{situation}[subsection]{Situation}

\theoremstyle{remark}
\newtheorem{remark}[subsection]{Remark}
\newtheorem{remarks}[subsection]{Remarks}

\numberwithin{equation}{subsection}

% Macros
%
\def\lim{\mathop{\mathrm{lim}}\nolimits}
\def\colim{\mathop{\mathrm{colim}}\nolimits}
\def\Spec{\mathop{\mathrm{Spec}}}
\def\Hom{\mathop{\mathrm{Hom}}\nolimits}
\def\Ext{\mathop{\mathrm{Ext}}\nolimits}
\def\SheafHom{\mathop{\mathcal{H}\!\mathit{om}}\nolimits}
\def\SheafExt{\mathop{\mathcal{E}\!\mathit{xt}}\nolimits}
\def\Sch{\mathit{Sch}}
\def\Mor{\mathop{\mathrm{Mor}}\nolimits}
\def\Ob{\mathop{\mathrm{Ob}}\nolimits}
\def\Sh{\mathop{\mathit{Sh}}\nolimits}
\def\NL{\mathop{N\!L}\nolimits}
\def\CH{\mathop{\mathrm{CH}}\nolimits}
\def\proetale{{pro\text{-}\acute{e}tale}}
\def\etale{{\acute{e}tale}}
\def\QCoh{\mathit{QCoh}}
\def\Ker{\mathop{\mathrm{Ker}}}
\def\Im{\mathop{\mathrm{Im}}}
\def\Coker{\mathop{\mathrm{Coker}}}
\def\Coim{\mathop{\mathrm{Coim}}}

% Boxtimes
%
\DeclareMathSymbol{\boxtimes}{\mathbin}{AMSa}{"02}

%
% Macros for moduli stacks/spaces
%
\def\QCohstack{\mathcal{QC}\!\mathit{oh}}
\def\Cohstack{\mathcal{C}\!\mathit{oh}}
\def\Spacesstack{\mathcal{S}\!\mathit{paces}}
\def\Quotfunctor{\mathrm{Quot}}
\def\Hilbfunctor{\mathrm{Hilb}}
\def\Curvesstack{\mathcal{C}\!\mathit{urves}}
\def\Polarizedstack{\mathcal{P}\!\mathit{olarized}}
\def\Complexesstack{\mathcal{C}\!\mathit{omplexes}}
% \Pic is the operator that assigns to X its picard group, usage \Pic(X)
% \Picardstack_{X/B} denotes the Picard stack of X over B
% \Picardfunctor_{X/B} denotes the Picard functor of X over B
\def\Pic{\mathop{\mathrm{Pic}}\nolimits}
\def\Picardstack{\mathcal{P}\!\mathit{ic}}
\def\Picardfunctor{\mathrm{Pic}}
\def\Deformationcategory{\mathcal{D}\!\mathit{ef}}
\usepackage{cleveref}
\crefname{originallemma}{original lemma}{original lemmas}
\Crefname{originallemma}{Original lemma}{Original lemmas}
\makeatletter
\newlabel{more-algebra-lemma-intersection-subfields-subspace}{{07P3}{source}{}{sourcepointer-1}{}}
\newlabel{more-algebra-lemma-intersection-subfields-subspace@cref}{{[originallemma][0][]07P3}{[1][0][]source}{}{}{}}
\makeatother
\title{A subspace of K to the n meets the base field if and only if it meets every subfield of a filtering family}
\author{The Stacks Project Authors}
\date{stacks-project:more-algebra.tex:local-snapshot-20261012; GNU FDL 1.2 or later}
\begin{document}
\maketitle
\setcounter{subsection}{1}
\noindent\textit{Editorial note.} All delivered mathematical text is the author's own, copied from the pinned Stacks Project chapter file more-algebra.tex; no statement or proof step was rewritten, reordered or removed.\par\medskip
\noindent\textit{Editorial note.} The chapter file pinned-original-source.tex carries only \textbackslash{}input\{preamble\} on its first line, so no preamble text lives inside the chapter. The retained file sources/stacks-more-algebra/retained-preamble.tex is a mechanical, line-for-line copy of the project's single shared preamble as staged in the snapshot, namely sources/stacks-chow/preamble.tex lines 14-18 (Xy-pic setup) and 154-217 (theorem declarations and mathematical macros). Nothing was added, removed or redefined, no package was added, and the excerpt serves only as preamble, never as a span coordinate.\par\medskip
\noindent\textit{Editorial note.} Both counted spans start at the \textbackslash{}begin line of their environment, so the statement and the author's proof are delivered whole; surrounding prose of the section is deliberately outside the counted spans and is carried as standard prerequisites instead.\par\medskip
\noindent\textit{Editorial note.} The citation fields (source\_id, source\_version, authors, source\_url, license.name, source\_title) are identical to those of the cards already published from this chapter, so that the batch selector treats them as one source instead of deferring them as a same-source clash.\par\medskip
\noindent\textit{Allowed standard prerequisites (editorial).}\begin{itemize}
\item The family of subfields is filtering downwards and has intersection k, with k contained in every member. hypotheses (1), (2) and (3) of the statement, source lines 11419-11424. Hypothesis (3) produces the smaller subfield K\_alpha' used in the last step, and (2) identifies the intersection of all subfields below K\_alpha with k.
\item Intersection with a subspace over a smaller field: if V cap K\_alpha to the n is spanned by a single vector with first coordinate 1, then elements of V over a subfield are scalar multiples of it. used at proof lines 11443-11447. Shows that V cap K\_alpha to the n equals K\_alpha v, which forces the later comparison with K\_alpha'.
\item Induction hypothesis applied to subspaces of K to the n-1 consisting of vectors whose first coordinate is zero. used at proof lines 11437-11440 to choose alpha with V cap 0 plus K\_alpha to the n-1 equal to zero. Produces the subfield that makes the first coordinate of the chosen vector nonzero.
\end{itemize}\medskip
\ifdefined\hypertarget\hypertarget{sourcepointer-1}{}\fi
\noindent\textit{Source reference.} Original lemma 07P3: focused-seeds/AGP/sources/stacks-snapshot-20261011/tags.txt, source lines 9783--9783. The lemma's own canonical tag, quoted so that readers can resolve the label inside the delivered text; it points at the counted statement itself and carries no proof content.\par\medskip
\begin{lemma}
\label{lemma-intersection-subfields-subspace}
Let $K/k$ be a field extension. Let $\{K_\alpha\}_{\alpha \in A}$
be a collection of subfields of $K$ with the following properties
\begin{enumerate}
\item $k \subset K_\alpha$ for all $\alpha \in A$,
\item $k = \bigcap_{\alpha \in A} K_\alpha$,
\item for $\alpha, \alpha' \in A$ there exists an $\alpha'' \in A$
such that $K_{\alpha''} \subset K_\alpha \cap K_{\alpha'}$.
\end{enumerate}
Then for $n \geq 1$ and $V \subset K^{\oplus n}$ a $K$-vector space
we have $V \cap k^{\oplus n} \not = 0$ if and only if
$V \cap K_\alpha^{\oplus n} \not = 0$ for all $\alpha \in A$.
\end{lemma}

\begin{proof}
By induction on $n$. The case $n = 1$ follows from the assumptions.
Assume the result proven for subspaces of $K^{\oplus n - 1}$.
Assume that $V \subset K^{\oplus n}$ has nonzero intersection with
$K_\alpha^{\oplus n}$ for all $\alpha \in A$. If
$V \cap 0 \oplus k^{\oplus n - 1}$ is nonzero then we win. Hence we may
assume this is not the case. By induction hypothesis we can find
an $\alpha$ such that $V \cap 0 \oplus K_\alpha^{\oplus n - 1}$
is zero. Let $v = (x_1, \ldots, x_n) \in V \cap K_\alpha^{\oplus n}$
be a nonzero element.
By our choice of $\alpha$ we see that $x_1$ is not zero.
Replace $v$ by $x_1^{-1}v$ so that $v = (1, x_2, \ldots, x_n)$.
Note that if $v' = (x_1', \ldots, x'_n) \in V \cap K_\alpha^{\oplus n}$, then
$v' - x_1'v = 0$ by our choice of $\alpha$. Hence we see that
$V \cap K_\alpha^{\oplus n} = K_\alpha v$. If we choose some
$\alpha'$ such that $K_{\alpha'} \subset K_\alpha$, then we
see that necessarily $v \in V \cap K_{\alpha'}^{\oplus n}$ (by the
same arguments applied to $\alpha'$). Hence
$$
x_2, \ldots, x_n \in
\bigcap\nolimits_{\alpha' \in A, K_{\alpha'} \subset K_\alpha} K_{\alpha'}
$$
which equals $k$ by (2) and (3).
\end{proof}

\end{document}
